\documentclass[11pt]{article}
\usepackage[T1]{fontenc}
\usepackage[utf8]{inputenc}
\usepackage{lmodern,microtype}
\usepackage{amsmath,amssymb,amsthm,mathtools}
\usepackage[margin=1in]{geometry}
\usepackage{booktabs,tabularx,array}
\usepackage{xcolor}
\usepackage[colorlinks=true,linkcolor=blue!45!black,citecolor=blue!45!black,urlcolor=blue!45!black]{hyperref}
\usepackage{enumitem}
\setlist{itemsep=3pt,topsep=5pt}
\setlength{\emergencystretch}{2em}
\newtheorem{theorem}{Theorem}[section]
\newtheorem{proposition}[theorem]{Proposition}
\theoremstyle{definition}
\newtheorem{definition}[theorem]{Definition}
\newtheorem{example}[theorem]{Example}
\DeclareMathOperator*{\argmax}{arg\,max}
\DeclareMathOperator{\pr}{pr}
\newcommand{\R}{\mathbb R}
\newcommand{\A}{\mathsf A}
\newcommand{\Cut}{\operatorname{Cut}}
\newcommand{\Eval}{\operatorname{Eval}}
\newcommand{\Opt}{\operatorname{Opt}}
\newcommand{\RootOpt}{\operatorname{RootOpt}}
\newcommand{\reporev}{777a28abb23ba9ccf9a98d0d14e0fb845907c5a8}
\newcommand{\reposrc}[2]{\href{https://github.com/eternal-november/eos-research-programmes/blob/\reporev/projects/contextual-composition/#1}{#2}}
\hypersetup{pdftitle={Prospective Interaction: Continuation Evaluation, Selection, and Native Composition},
 pdfauthor={},pdfsubject={A finite-history pilot for proactads}}
\title{Prospective Interaction\\
\large Continuation Evaluation, Selection, and Native Composition}
\author{}
\date{Working mathematical note --- 26 November 2026}
\begin{document}
\maketitle
\vspace{-1.5em}
\begin{abstract}
Anticipatory control requires a model of possible continuations and a rule
for choosing actions using their consequences. A scalar value calculation
can compose even when the corresponding policy interface loses information:
a three-stage example has four policies in one bracketing and two in the
other, although both give the same optimal value for every terminal score.
This note develops a finite pilot that retains complete policy tables on an
explicitly authorized history interface. Evaluation composes across complete
cuts; all-history optimal policies form a wide subcategory of the category
of evaluated policies. An exact loss formula separates root optimality from
optimality at every retained history. A sharp rectangularity criterion
identifies when local choices can be assembled, and a value-comparison
theorem for independently supplied native maps identifies control-coverage and policy-descent
conditions. The Bellman recursion is standard; the purpose is to delimit its
use in the contextual-composition programme. These results support prospective
equipment on a finite controlled realization. A lifted context/content
presentation, general feedback, and model or reference revision remain
separate obligations.
\end{abstract}

\section{Scope, evidence, and the proposed mathematical interface}
\label{sec:scope}
The design question is whether prediction and anticipatory regulation require
a broader interaction calculus, or additional structure on an existing one.
The present pilot makes that question testable at a restricted interface.
Its input is a finite, fully retained controlled history arena, a fixed
predictive row for each complete control, edge rewards, and a terminal
continuation valuation. Its output includes evaluated \emph{complete policies},
their optimal subfamily, and comparison laws. The scalar value alone is an
insufficient output: it can hide a failure to preserve implementable policies.

Capucci and Myers construct contextful arrows through contextads and wreath
products \cite{contextads}. The contextual-composition programme's
provisional interactad package retains
one-sided context and content presentations and a full pseudodistributive law
\(\tau:N\odot_C M\Rightarrow M\odot_C N\), including native cells and the unit
and multiplication coherence data \cite{terminology,relative}. That interface
sets the eventual standard for a proactad construction. A Bellman equation
does not itself supply such a crossing. The construction below is deliberately
one level earlier: it establishes what finite prospective evaluation and
selection can contribute to a supplied controlled realization.

\paragraph{Pinned assumptions.}
Every retained decision history has a finite nonempty menu of complete
controls, including any admitted law label. The structural history tree is
finite. Raw zero-probability outcomes and histories behind unchosen controls
remain distinct. Current revealed history is explicitly authorized as the
controller's observation. Controls select fresh-outcome rows; an existing
reserve is never redrawn. Where a row represents a joint draw, its central
joint service is an explicitly supplied interface. Rewards and continuation
values are fixed real-valued inputs.

These assumptions restrict both information and assembly. Merely retaining
an audit history does not authorize a controller to read it. Finitely many
raw outcomes do not imply finitely many admitted probability-law labels.
If the supplied rows are forecasts, all optimality statements concern those
forecasts; equality with the actual execution law is an additional condition.
The scalar reward is a supplied evaluation criterion, not a proof that a
stakeholder's substantive requirement is satisfied. Hard admissibility must
be encoded separately in the legal control doctrine.

\paragraph{Declared kernel and residuals.}
The pilot holds the primitive network, observation permission, predictive
model, and objective fixed. It does not construct coefficient or grade maps,
infer a controller's authority, identify observationally indistinguishable
histories, or learn a model. It proves composition across cuts of \emph{one
common arena}; it does not compose arbitrary opaque component controllers.
Infinite horizons, arbitrary whole-plan commitments, decentralized
information, and reference revision require further interfaces.

\paragraph{Evidence classes.}
The propositions and theorems below have written finite proofs. The companion
Python program supplies exact rational diagnostics for specified examples,
including independent path enumeration against backward recursion. No new
proof-assistant verification is claimed. The source-reading route and its
pinned revision are supplied separately.

\begin{center}
\small
\begin{tabularx}{\textwidth}{@{}lXl@{}}
\toprule
ID & Question answered & Location\\
\midrule
PI-0 & Can equal optimized values hide failed policy associativity? & \S\ref{sec:obstruction}\\
PI-1 & What composes across a retained cut? & \S\ref{sec:cuts}\\
PI-2 & How do root and all-history selection differ? & \S\ref{sec:selection}\\
PI-3 & Which actual categories does the pilot supply? & \S\ref{sec:categories}\\
PI-4 & When is independent local maximization exact? & \S\ref{sec:rectangle}\\
PI-5 & What survives a supplied native comparison? & \S\ref{sec:native}\\
PI-6 & When does an optimal source policy extend or descend? & \S\ref{sec:fibre}\\
\bottomrule
\end{tabularx}
\end{center}

\section{A value-preserving failure of policy associativity}
\label{sec:obstruction}
Consider the tempting substitution rule for a controlled process
\(f:X\to Y\) followed by \(g:Y\to Z\):
\begin{equation}
 A_{g\circ f}(x)
   =\coprod_{a\in A_f(x)}\prod_{y\in Y} A_g(y).
 \label{eq:naive}
\end{equation}
It records an initial control and a continuation control for each
intermediate output. For sequential assembly, the question is whether the
intermediate output retains the information needed by subsequent controls.

\begin{proposition}[Naive substitution obstruction; PI-0]
\label{prop:naive}
Rule \eqref{eq:naive}, with only the displayed external interfaces, need
not admit a policy-preserving associator. This failure can remain invisible
to every optimized terminal-value calculation.
\end{proposition}
\begin{proof}
Let \(f:1\to\{0,1\}\) have one control and emit a fair bit. Let
\(g:\{0,1\}\to1\) have one control and forget the bit. Let
\(h:1\to\{0,1\}\) have two controls, deterministically emitting \(0\) and
\(1\). All rewards are zero.

In \((h\circ g)\circ f\), a continuation control of \(h\circ g\) is chosen
for each output of \(f\). There are four tables
\[
 (0,0),\quad(0,1),\quad(1,0),\quad(1,1).
\]
In \(h\circ(g\circ f)\), the interface before \(h\) is a singleton. There
are only two constant tables, \((0,0)\) and \((1,1)\). Thus no bijection
of the retained policy sets exists. Their attainable output laws also
differ: the first bracketing permits probabilities
\(\{0,\tfrac12,1\}\) of output \(1\), while the second permits
\(\{0,1\}\).

For every \(v:\{0,1\}\to\R\), however, each policy's expected value is a
convex combination of \(v(0)\) and \(v(1)\). Both bracketings contain both
constant policies, so both optimal values equal
\(\max\{v(0),v(1)\}\).
\end{proof}

This is an obstruction to the proposed interface, not an impossibility
theorem for prospective interaction. A common retained-history arena, or
an explicitly typed controller-memory boundary, can preserve the relevant
bit across both bracketings. If the controller is forbidden to observe or
retain that bit, rebracketing cannot grant permission: the admissible table
family must enforce the restriction. Section~\ref{sec:rectangle} measures
the resulting assembly condition.

\section{Finite retained arenas and exact cut evaluation}
\label{sec:cuts}
\begin{definition}[Finite controlled history arena]
\label{def:arena}
A finite controlled history arena \(\A\) is a finite rooted tree (or
disjoint forest) with decision histories \(H^{\mathrm{dec}}\) and terminal
histories \(L\). At each decision history \(h\), there is a finite nonempty
set \(A(h)\) of complete controls. For \(a\in A(h)\), there are:
\[
 W(h,a),\qquad
 q_{h,a}:W(h,a)\to[0,1],\qquad
 \sum_{w\in W(h,a)}q_{h,a}(w)=1,\qquad
 r(h,a,w)\in\R.
\]
The finite raw carrier \(W(h,a)\) is retained in full. Its successor
\(h[a,w]\) records the full history extension, including the control label
and raw witness. Different such extensions are distinct tree nodes.
Terminal histories have no controls and retain their declared completion
or forced-stop status. A terminal valuation is \(v:L\to\R\).
\end{definition}

A deterministic complete policy is an element
\[
 \Pi_\A=\prod_{h\in H^{\mathrm{dec}}}A(h).
\]
Every coordinate is present before any policy is selected. A zero history
has zero mass under a particular selected row; an off-policy history lies
behind a control not selected by that policy. Neither distinction deletes
a coordinate. Probabilistic service rows may themselves be controls in the
finite menu; randomization over policy tables is not added implicitly.

For a supplied policy \(\pi\), evaluation and optimal evaluation are defined
by backward recursion:
\begin{align}
 V^\pi(h)&=v(h), & V^*(h)&=v(h) &&(h\in L),\notag\\
 V^\pi(h)&=\sum_w q_{h,\pi(h)}(w)
       \bigl(r(h,\pi(h),w)+V^\pi(h[\pi(h),w])\bigr), \label{eq:eval}\\
 Q(h,a)&=\sum_w q_{h,a}(w)
       \bigl(r(h,a,w)+V^*(h[a,w])\bigr),
 &V^*(h)&=\max_{a\in A(h)}Q(h,a). \label{eq:bellman}
\end{align}
Finite well-foundedness makes both recursions defined. No stationarity,
discounting, or infinite-horizon convergence is being assumed.

A \emph{complete cut} \(K\) is an antichain meeting every structural
root-to-terminal path exactly once, including paths through zero outcomes
and through unchosen controls. Let \(F\) be the prefix stopped immediately
before acting at \(K\), and let \(G_k\) be the subtree beginning at \(k\).
If \(k\) is terminal, \(G_k\) has the unique empty policy.

\begin{theorem}[Retained-cut composition; PI-1]
\label{thm:cut}
There is a canonical bijection
\begin{equation}
 \Pi_\A\ \cong\ \Pi_F\times\prod_{k\in K}\Pi_{G_k}.
 \label{eq:policycut}
\end{equation}
For a root \(h_0\), write \(R_F(\pi_F)\) for expected reward before the
cut and \(p_{\pi_F}(k)\) for its hitting law. Then
\begin{equation}
 V^\pi(h_0)=R_F(\pi_F)+
       \sum_{k\in K}p_{\pi_F}(k)V^{\pi_k}_{G_k}(k).
 \label{eq:cutvalue}
\end{equation}
Consequently, if \(B\) denotes optimal continuation evaluation,
\begin{equation}
 B_\A(v)=B_F\bigl((B_{G_k}(v|_{L(G_k)}))_{k\in K}\bigr).
 \label{eq:bellmancut}
\end{equation}
For nested cuts, the policy bijections obey associativity and unit
coherence while retaining all coordinates.
\end{theorem}
\begin{proof}
The decision histories partition into those strictly before \(K\) and
those in the subtrees rooted at \(K\). Restriction of a policy to these
disjoint coordinate sets gives \eqref{eq:policycut}; union is its inverse.
Every selected raw path has a unique first meeting with \(K\). Factoring
its product probability at that meeting and splitting its accumulated
reward proves \eqref{eq:cutvalue}, including terms of mass zero.

For each \(k\), backward recursion gives a policy attaining the largest
continuation value at every history of \(G_k\): choose a maximizing control
at each node. These choices can be made simultaneously because the policy
coordinates form the displayed product. Substituting their values into
\eqref{eq:cutvalue} and maximizing the prefix proves
\eqref{eq:bellmancut}. Conversely, no supplied suffix exceeds its optimal
value, so the same expression bounds every complete policy.

For nested cuts, either order of restriction records the same controls at
the same primitive histories. Union has the same property. All
associativity and unit diagrams therefore commute as coordinate identities.
\end{proof}

This is standard finite dynamic programming with the policy coordinates
kept visible. Its assembly input is a common primitive-history arena. The
project's all-policy open-assembly theorem likewise concerns a presented
primitive network and its full policy family, rather than an operation on
arbitrary isolated opaque operand policies \cite{transport}.

\section{Root optimality and optimality at every history}
\label{sec:selection}
Define the local greedy sets and nonnegative losses by
\[
 G(h)=\argmax_{a\in A(h)}Q(h,a),\qquad
 \delta(h,a)=V^*(h)-Q(h,a)\geq0.
\]
A policy is \emph{all-history greedy} when \(\pi(h)\in G(h)\) for every
decision history. This is equivalent to attaining \(V^*(h)\) from every
history \(h\), including histories of probability zero under the root
execution. A \emph{root-optimal} policy is only required to attain
\(V^*(h_0)\) at the specified root.

\begin{theorem}[Exact policy loss and selection fibres; PI-2]
\label{thm:loss}
Let \(d_\pi(h)\) be the probability that execution from \(h_0\) reaches
\(h\). For every complete policy,
\begin{equation}
 V^*(h_0)-V^\pi(h_0)
   =\sum_{h\in H^{\mathrm{dec}}}d_\pi(h)\delta(h,\pi(h)).
 \label{eq:loss}
\end{equation}
Hence root optimality is equivalent to greediness at every history with
\(d_\pi(h)>0\). The all-history greedy family is
\[
 G_\A(v)=\prod_{h\in H^{\mathrm{dec}}}G(h).
\]
For the cut in Theorem~\ref{thm:cut}, write \(V_G(k)=B_{G_k}(v)\),
with \(v\) restricted to the relevant subtree's terminals.
Suffix root optimality is relative to each suffix's own root.
Then
\begin{equation}
 G_\A(v)\cong G_F(V_G)\times\prod_{k\in K}G_{G_k}(v).
 \label{eq:greedycut}
\end{equation}
The root-optimal family instead has the exact fibre decomposition
\begin{equation}
\begin{split}
 \RootOpt_\A(v)\cong
 \coprod_{\pi_F\in\RootOpt_F(V_G)}
 \left(
 \prod_{k:\,p_{\pi_F}(k)>0}\RootOpt_{G_k}(v)
 \ \times\
 \prod_{k:\,p_{\pi_F}(k)=0}\Pi_{G_k}
 \right).
\end{split}
\label{eq:rootfibre}
\end{equation}
\end{theorem}
\begin{proof}
Set \(\Delta_\pi(h)=V^*(h)-V^\pi(h)\). Subtracting the two evaluations at
the selected control gives
\[
 \Delta_\pi(h)=\delta(h,\pi(h))+
       \sum_wq_{h,\pi(h)}(w)\Delta_\pi(h[\pi(h),w]).
\]
At terminals \(\Delta_\pi=0\). Finite expansion gives
\eqref{eq:loss}; every summand is nonnegative. Equality at the root
therefore has exactly the stated positive-occupancy criterion.

All-history greediness restricts to precisely the local greedy conditions
in each disjoint segment; suffix optimal values are the continuation values
used by the prefix. This proves \eqref{eq:greedycut}. For
\eqref{eq:rootfibre}, subtract \eqref{eq:cutvalue} from the optimum:
\[
\begin{split}
 V^*_\A(h_0)-V^\pi_\A(h_0)
 ={}&\left[B_F(V_G)-E_F^{\pi_F}(V_G)\right](h_0)\\
 &+\sum_kp_{\pi_F}(k)
       \left[V_G(k)-V_{G_k}^{\pi_k}(k)\right].
\end{split}
\]
Here \(E_F^{\pi_F}\) is evaluation with the prefix policy fixed. The
bracketed terms are nonnegative. Their weighted sum vanishes exactly
when the prefix is root-optimal and every positive-mass suffix is
root-optimal. Zero-mass suffix policies remain unrestricted.
\end{proof}

All-history selection is a substantive requirement. It can express
contingency-wise preparedness or support later re-entry at any retained
history. It cannot be silently substituted for the weaker root objective.
In the companion depth-three arena, the complete table set has \(8192\)
members, \(1024\) are root-optimal, and exactly one is all-history greedy.
The retained zero and off-policy coordinates account for the difference.

\section{Categories of evaluated and selected policies}
\label{sec:categories}
Fix the entire arena \(\A\). Its complete cuts form a finite poset:
\(C\leq D\) when every structural path meets \(C\) no later than \(D\).
Let \(H_{CD}\) contain the decision histories from \(C\) inclusive to
\(D\) exclusive; execution stops before acting at \(D\). Let
\(\Pi_{CD}=\prod_{h\in H_{CD}}A(h)\). A shared terminal node simply
persists across the cuts, with an empty segment.

For \(\pi\in\Pi_{CD}\), evaluation defines a monotone affine map
\[
 E_{CD}^{\pi}:\R^D\longrightarrow\R^C.
\]
Optimal evaluation \(B_{CD}:\R^D\to\R^C\) is monotone. The finite cut
theorem gives, for \(C\leq D\leq E\),
\begin{align}
 E_{CE}^{\pi\amalg\rho}
     &=E_{CD}^{\pi}\circ E_{DE}^{\rho},
       &B_{CE}&=B_{CD}\circ B_{DE}, \label{eq:segmentcomp}\\
 G_{CE}(w)&\cong
       G_{CD}(B_{DE}w)\times G_{DE}(w).
       \label{eq:segmentgreedy}
\end{align}
The symbol \(\amalg\) here means union of tables on disjoint history
coordinates. In particular \(C\mapsto\R^C\), with arrows \(B_{CD}\),
is a functor \(\Cut(\A)^{\mathrm{op}}\to\mathbf{Pos}\).

\begin{theorem}[Evaluation category and optimal subcategory; PI-3]
\label{thm:categories}
There is a category \(\Eval(\A)\) whose objects are pairs \((C,u)\) with
\(u\in\R^C\), and whose arrows
\[
 (C,u)\longrightarrow(D,v)
\]
are policies \(\pi\in\Pi_{CD}\), for \(C\leq D\), satisfying
\(u=E_{CD}^{\pi}(v)\). Composition is table union.

There is a wide, generally non-full subcategory \(\Opt(\A)\) in which an
arrow is required to satisfy
\[
 \pi\in G_{CD}(v),\qquad u=B_{CD}(v).
\]
For a fixed optimal arrow from \((C,u)\) to \((E,w)\) and an intermediate cut
\(D\), its factors are uniquely its restrictions, through the object
\((D,B_{DE}(w))\).
\end{theorem}
\begin{proof}
Equation~\eqref{eq:segmentcomp} makes composition well typed in
\(\Eval(\A)\). Disjoint table union is associative; the unique empty
table at \(C=D\) evaluates by the identity and supplies the identity
arrow at \((C,u)\).

An all-history greedy policy evaluates to \(B_{CD}(v)\), so every
\(\Opt(\A)\) arrow is an \(\Eval(\A)\) arrow. Equation
\eqref{eq:segmentgreedy} proves closure under composition and the stated
factorization. Empty tables are greedy, so the subcategory is wide.
For non-fullness, take a single-root arena with a root-optimal table
having a nongreedy off-policy coordinate: it has the same endpoint
valuations as a greedy table while remaining outside \(\Opt(\A)\).
\end{proof}

This category retains supplied, possibly suboptimal policies. Optimization
is additional selected structure. At a fixed cut \(C\), an identity-segment
arrow exists only when \(u=v\). Thus the construction does \emph{not} already
contain same-time changes of valuation, reference, or model. It also makes
no fibration assertion: a factorization of an already supplied complete
table is different from a universal lift of an arbitrary external map.
The arena parameter remains fixed throughout.

\section{A sharp criterion for assembling local choices}
\label{sec:rectangle}
Information restrictions, commitments, or coupled feasibility may replace a
product of local policy sets by a proper subset. A scalar Bellman step then
solves a relaxation unless the required combination of local choices is
actually admissible.

\begin{theorem}[Positive-coordinate rectangularity; PI-4]
\label{thm:rectangle}
Let \(K\) be finite, \(p\) a probability law on \(K\), each \(U_k\) a finite
nonempty set of local alternatives, and
\(\varnothing\ne S\subseteq\prod_{k\in K}U_k\) the admissible complete
table family. For scores \(u_k:U_k\to\R\), let
\[
 J_S(u)=\max_{s\in S}\sum_kp_k u_k(s_k),\qquad
 J_{\mathrm{sep}}(u)=\sum_kp_k\max_{a\in U_k}u_k(a).
\]
Then \(J_S(u)\leq J_{\mathrm{sep}}(u)\). For a fixed score family,
equality holds precisely when \(S\) contains a table selecting a local
maximizer at every positive-mass coordinate. In that case its maximizing
tables are exactly the members of \(S\) with this property.

Write \(K_+=\{k:p_k>0\}\). Equality for \emph{every independent family
of real score tables} is equivalent to
\begin{equation}
 \pr_{K_+}(S)=\prod_{k\in K_+}U_k.
 \label{eq:rectangle}
\end{equation}
\end{theorem}
\begin{proof}
For a supplied \(s\), the gap from the separated value is
\[
 \sum_k p_k\bigl(\max_a u_k(a)-u_k(s_k)\bigr)\geq0.
\]
It vanishes exactly at positive-coordinate local maximizers. Finite
nonemptiness makes the maximum over \(S\) attained and proves the
fixed-score statement.

If \eqref{eq:rectangle} holds, choose a local maximizer at every
positive coordinate and extend that tuple to \(S\). This gives equality.
Conversely, if a tuple \(a^*\) is missing from the positive-coordinate
projection, take \(u_k(a)=\mathbf1_{\{a=a^*_k\}}\) there and zero scores
elsewhere. The separated value is \(1\). Every admissible table misses
at least one positive coordinate, giving value at most
\(1-\min_{k\in K_+}p_k<1\). Uniform equality fails.
\end{proof}

\begin{example}[Hidden observation or prior commitment]
\label{ex:diagonal}
Let \(p=(\tfrac12,\tfrac12)\), \(U_0=U_1=\{0,1\}\), and
\(S=\{(0,0),(1,1)\}\). With \(u_0(b)=1-b\) and \(u_1(b)=b\),
the constrained value is \(1/2\) while independent maximization gives
\(1\), using the forbidden table \((0,1)\). The same admissible family
can arise from different mechanisms: an unobserved bit, a prior commitment
to one common action, or an explicit cross-history constraint. A repair
must preserve the actual mechanism. Observation permission can change the
first problem; it need not relax either of the other two.

For \(p=(1,0)\), the same \(S\) satisfies \eqref{eq:rectangle}.
Its scalar evaluation is exact for all scores even though it contains only
two of the four retained raw tables. Value equality therefore does not
certify full native table coverage.
\end{example}

The necessity result concerns arbitrary independent local scores. A fixed
model with only a restricted family of terminal valuations may fail to
distinguish omitted duplicate or dominated alternatives. Likewise, if
feasibility depends on the prefix policy, the theorem applies to each
fixed prefix and its actual suffix family; it does not remove that
dependency.

\section{Evaluation along independently supplied native maps}
\label{sec:native}
The project already distinguishes primitive/native comparison data from
the scalar conditional-law equations that those data must satisfy
\cite{conditional,relative}. The following theorem evaluates a supplied
comparison. It does not reconstruct its coefficient, full-grade, world,
or boundary arrows from probability rows.

Let \(S,T\) be finite arenas. Suppose a history map \(F:H_S\to H_T\)
preserves depth and terminal/forced-stop status, and for each source
decision history there are maps
\[
 \phi_h:A_S(h)\to A_T(Fh),\qquad
 \rho_{h,a}:W_S(h,a)\to W_T(Fh,\phi_h a).
\]
Require exact compatibility:
\begin{align}
 (\rho_{h,a})_*q^S_{h,a}&=q^T_{Fh,\phi_h a},
       \label{eq:push}\\
 F(h[a,w])&=(Fh)[\phi_h a,\rho_{h,a}w], \label{eq:body}\\
 r_S(h,a,w)&=r_T(Fh,\phi_h a,\rho_{h,a}w),\qquad
 v_S=v_T\circ F\quad\hbox{on terminals}. \label{eq:reward}
\end{align}
For a \emph{native} interpretation, the underlying comparison in the
chosen source doctrine, with its material and interface squares, is an
additional supplied hypothesis.

\begin{theorem}[Control coverage and value comparison; PI-5]
\label{thm:native}
Under \eqref{eq:push}--\eqref{eq:reward},
\[
 V^*_S(h)\leq V^*_T(Fh).
\]
If every \(\phi_h\) is surjective, then at every source history and control,
\begin{equation}
 V^*_S(h)=V^*_T(Fh),\qquad
 Q_S(h,a)=Q_T(Fh,\phi_h a),\qquad
 G_S(h)=\phi_h^{-1}(G_T(Fh)).
 \label{eq:nativeequal}
\end{equation}
No raw-map surjectivity is required for these scalar statements. These
evaluated comparison data, and their control-surjective subclass, are
closed under identities and composition.

For already compatible policies \(\pi_S,\pi_T\), meaning
\(\pi_T(Fh)=\phi_h(\pi_S(h))\) at every source decision history,
fixed-policy values agree:
\[
 V^{\pi_S}_S(h)=V^{\pi_T}_T(Fh),
\]
even without control surjectivity.
\end{theorem}
\begin{proof}
Induct backward on the source tree. Terminal values agree by hypothesis.
For a source control, the induction inequality at its successors, followed
by exact pushforward and reward preservation, gives
\[
 Q_S(h,a)\leq Q_T(Fh,\phi_h a).
\]
Maximizing over source controls gives the stated value inequality.
If every control map is onto, the induction starts with equality at
terminals and gives equality for every source control. Surjectivity then
identifies the two maxima and their maximizing preimages, proving
\eqref{eq:nativeequal}.

A target raw atom outside the image of \(\rho_{h,a}\) has zero mass by
\eqref{eq:push}; it contributes nothing to this scalar sum. It remains an
actual target history. Identity and composite maps satisfy the same
equations because pushforwards and the successor and reward squares
compose; composites of surjective control maps are surjective.

For compatible fixed policies, the induction compares the selected rows
only, with no maximization. It yields equality without any coverage
condition on unselected controls.
\end{proof}

Control coverage is a sufficient structural condition, not a necessary
condition for scalar equality. A target may add a duplicate control or a
strictly dominated control without changing any optimal value. At a
one-stage boundary each control contributes an affine profile
\(L_a(v)=\bar r_a+p_a\cdot v\); the Bellman operator sees only their upper
envelope. Raw labels, grades, and dominated or duplicated alternatives
can survive outside that observation.

\subsection{The exact greedy extension fibre}
\label{sec:fibre}
An equality of optimal values does not choose a compatible target policy.
For a supplied complete source policy \(p\), define the forced controls at
a target history by
\[
 D_k(p)=\{\phi_h(p(h)):Fh=k\}.
\]
Let \(G_T(k)\) use the fixed target model and continuation valuation.

\begin{theorem}[Greedy extension and descent; PI-6]
\label{thm:fibre}
The full family of all-history greedy target policies compatible with \(p\)
is exactly
\begin{equation}
 \prod_{k\in H_T^{\mathrm{dec}}}
 \left(
 G_T(k)\cap
 \{b\in A_T(k): b=d\text{ for every }d\in D_k(p)\}
 \right).
 \label{eq:extfibre}
\end{equation}
If the control maps are surjective and \(p\) is all-history greedy, this
family is nonempty if and only if the forced controls agree on every
nonempty \(F\)-fibre. Covered target coordinates are then forced.
Uncovered coordinates retain their full greedy sets. When an extension
exists, it is unique if and only if each such set is a singleton.
\end{theorem}
\begin{proof}
Compatibility gives exactly the displayed constraints at each target
history. The target's complete policy space is a product, so the constraints
assemble as \eqref{eq:extfibre}. Under the additional assumptions,
Theorem~\ref{thm:native} puts every forced image in \(G_T(k)\). A covered
factor is therefore nonempty precisely when all its forced values agree;
then it is a singleton. An uncovered factor is \(G_T(k)\), which is finite
and nonempty. The existence and uniqueness statements follow.
\end{proof}

This is an optimizer-constrained finite analogue of the existing relative
extension-fibre theorem \cite{conditional}. That source theorem fixes the
target structural skeleton; the present control coordinates can bundle
structural choices and law labels. Identifying those two presentations
requires the appropriate normal-form interpretation.
Surjectivity on histories removes uncovered-coordinate freedom but
does not remove disagreement between merged source histories.

\subsection{Optimization does not repair a missing intermediate policy}
\label{sec:middle}
The source policy-transport relation satisfies
\(R_gR_f\subseteq R_{gf}\); equality requires an intermediate complete
policy \cite{transport}. A simple optimal version of the source's strict
witness survives the new selection condition.

After a fair root bit, a source policy chooses at each of two histories a
law on four raw bits \((X,Y,U,V)\). Each choice is a product
\(L\otimes H\), where \(L\in\{D,A\}\) couples \((X,Y)\) and
\(H\in\{\mathrm{good},\mathrm{bad}\}\) couples \((U,V)\).
The first two couplings are
\[
 D(0,0)=D(1,1)=\tfrac12,\qquad
 A(0,1)=A(1,0)=\tfrac12,
\]
with the other two raw weights zero. Good is the diagonal coupling on
\((U,V)\); bad is the antidiagonal coupling. All sixteen raw source
atoms are retained, including zeros. Every bit marginal is fair.
Root reward and terminal continuation values are zero; the second-stage
reward is the pointwise function \(\mathbf1_{\{U=V\}}\). Its expectation
is \(1\) under good and \(0\) under bad, independently of \(L\).
The first comparison merges the two root histories but retains the joint
four-bit outcomes. A middle policy must consequently choose the same full
law at both source histories. The second comparison projects
\((X,Y,U,V)\mapsto(U,V)\), forgetting \(L\) while retaining the actual
law \(H\). This projection preserves the reward at every raw atom,
including atoms of zero mass. The structural skeleton remains fixed;
good and bad are admitted laws, not extra reward tags. All material
coefficients in this witness may be units.

The source policy
\[
 ((D,\mathrm{good}),(A,\mathrm{good}))
\]
and target policy \(\mathrm{good}\) are compatible under the composite.
Both are all-history optimal. No middle policy is compatible with the
source policy, because it would need both distinct coupling labels.
The finite diagnostic finds four compatible optimal source--target pairs,
of which only two factor through the middle. The nonconstant reward makes
selection nonvacuous. Thus restricting to optimal policies does not make
the existing policy relation compose strictly. A general policy-lifting
or fibration claim would need additional hypotheses.

\section{What the finite result does not identify}
\label{sec:boundaries}
\subsection{Fresh execution versus a whole-plan commitment}
The fresh-row convention in Definition~\ref{def:arena} matters. If a label
instead carries a whole joint law on an already retained reserve \(R\)
and fresh witnesses \(W\), execution must retain the actual reserve and
sample the conditional fresh law. Assigning the whole law to the fresh
service kernel would redraw the reserve.

The source nonblocking doctrine supplies calibration equation C1 and the
gluing identity C2, along with incoming-signature validity and pointwise transport of fresh
service tables, including zero tables \cite{nonblocking}. Authorizing the
identity history view does not establish these conditions. A Bellman
maximizer among individually admitted whole plans need not be a calibrated
selector of those plans. Applying this pilot to that doctrine requires
an interpretation proving calibration and admissible assembly for the
selected policy family.

\subsection{Finite cut evaluation versus horizon revision}
At time \(0\), setup costs \(1\); at time \(1\), completion pays \(3\)
only if setup occurred. Waiting pays zero, and there is no terminal
salvage value. The two-step optimum is setup then completion, returning
\(2\). Repeated one-step optimization with terminal valuation zero
chooses waiting at both times, returning zero.

This does not contradict the cut theorem: the correct continuation
values at time \(1\) are \(3\) for a ready state and \(0\) for an unready
state. Resetting both to zero changes the optimization problem.
Horizon-one and horizon-two optima need not even agree at the root.
Consequently, independent finite-horizon optimization does not supply
the coherent controller family required by the source's native
infinite-prefix/finality result \cite{nonblocking}.

\subsection{Concurrent predictions need their joint law}
Let \(D_1,D_2\) be binary demands, each with marginal probability \(1/2\).
Capacity \(c\in\{1,2\}\) has cost
\[
 J(c)=\tfrac34c+2\Pr(D_1+D_2>c).
\]
Independent demands give
\(J(1)=5/4<J(2)=3/2\), whereas a common shock \(D_1=D_2\) gives
\(J(1)=7/4>J(2)=3/2\). The component marginals coincide and the
optimal joint decision reverses.

Thus a prospective concurrency interface must retain the joint law or
other data sufficient for the actual joint objective. With a genuinely
factorized feasible set and an additive sum of separate local objectives,
the optimizer factors as a product of local optimizer sets; the proof is
termwise maximization. Shared constraints or a nonseparable cost remove
that premise. This example is a decision-theoretic obstruction, not a
construction of the programme's concurrency tensor.

\subsection{Prediction, regulation, and revision are distinct operations}
Feedback specifies how interaction recurs. Prospective evaluation uses
consequences beyond the present boundary. Regulation additionally identifies
what is to be maintained or prepared. Reference or model revision changes
those supplied data. None of these operations follows merely from naming
another one.

For this design study, allostatic equipment means a proposed interface for
context-sensitive anticipatory regulation; reference revision is a possible
additional operation, not a physiological theorem or a prerequisite imposed
on every anticipatory controller. The setup example already models
preparation for a later benefit with a fixed objective. To add changing
references or predictive models, the structure must state who or what
authorizes the change, what observation triggers it, how continuation
valuations are transported, and how existing commitments remain valid.
The same-cut boundary of \(\Eval(\A)\) exhibits missing valuation-revision
arrows. Model revision changes the arena parameter itself and requires
comparisons between differently modelled arenas.

\section{Consequence for a prospective interaction programme}
\label{sec:consequence}
The finite pilot supports a \emph{prospectively equipped controlled
realization}: retained admissible controls and predictions, continuation
evaluation, and a separately specified selection relation. It supports
the specialized-structure branch of the proactad question. Adding such
equipment supplies more structure and hypotheses on an interaction system;
that fact alone does not exhibit a more general class of underlying
interaction systems.

The distinction is testable. On a fixed arena with zero rewards and zero
terminal values, every policy is greedy, so the selected zero-continuation
fibre reduces to the full policy family on that arena. This is a degenerate reduction
inside the common finite host. It is not an embedding theorem for arbitrary
reactads or interactads, since finiteness, observation permission, and
fresh-row assembly remain assumptions.

Categorical cybernetics offers an established comparison point:
parametrized optics organize systems with choices and selection relations
\cite{cybernetics}; work on categorical reinforcement learning treats
Bellman operators through optics and a contravariant representable
construction \cite{rl}. The present finite recursion is not a novelty
claim relative to that literature. The question specific to this programme
is whether that evaluation/selection structure can preserve the richer
native presentations, cells, law labels, and controlled forgetting.
Proposition~\ref{prop:naive} and Theorem~\ref{thm:fibre} show why value
semantics alone cannot answer it.

A full proactad proposal in the programme's sense would still need the
following concrete extension of the interface:
\begin{enumerate}[label=(\roman*)]
\item \textbf{Presentation and forgetful comparison.}
Specify candidate prospective context and content presentations, their
pure endpoint actions, and a forgetful comparison to the chosen interaction
doctrine. Declare which predictions, objectives, and selection witnesses
are objects or cells, and which are external parameters.
\item \textbf{Crossing and coherence.}
Construct the proposed crossing on native cells, with pure-middle balance
and all unit/multiplication modifications. A scalar Bellman identity is
only a required evaluation equation for such a construction.
\item \textbf{Admissible selection and transport.}
Prove the relevant analogue of rectangularity, or retain the constrained
selection relation when it fails. Prove calibration for committed whole
plans and account for merged and uncovered histories using actual lift
fibres.
\item \textbf{Feedback and optional revision.}
Specify a consistent finite or infinite continuation doctrine. If model
or reference revision is included, specify its authorized arrows and the
preservation obligations on existing commitments. State stability, viability, or regret
claims only with their own hypotheses and evidence.
\end{enumerate}

The next discriminating extension is therefore a presentation-level
interpretation of this finite category of policies and its selector,
with explicit observation and commitment interfaces. If that
interpretation closes under the existing native crossing, prospective
interaction is additional equipment. If it fails, the failed square,
lift fibre, or coherence law can identify precisely which generalization
is required. The current obstructions justify preserving more interface
data; they do not yet require replacing the existing interaction calculus.

\appendix
\section{Executable diagnostics and reproducibility}
\label{sec:checks}
The companion \texttt{checks/check\_finite\_prospective.py} uses only the
Python standard library and \texttt{fractions.Fraction}. Run:
\begin{quote}
\texttt{python3 checks/check\_finite\_prospective.py}
\end{quote}
It writes \texttt{checks/results.json}. Ordinary Python execution is
required; optimized mode that disables assertions is inappropriate for
these diagnostics.

\begin{center}
\small
\begin{tabularx}{\textwidth}{@{}lX@{}}
\toprule
Diagnostic & Executed result\\
\midrule
Full retained arena & \(13\) decision histories, \(16\) terminal histories,
\(8192\) complete policies; direct enumeration and backward evaluation
both give value \(2\).\\
Loss and cut coordinates & All \(8192\) loss identities and table
split/reassembly checks pass; \(1024\) root optima, one all-history optimum.\\
Rectangularity & All \(15\) nonempty binary table families, under three
prefix laws; \(720\) score comparisons.\\
Associativity obstruction & Four versus two policies; attainable final
probabilities \(\{0,\tfrac12,1\}\) versus \(\{0,1\}\).\\
Native comparison & Strict comparison without control coverage;
duplicate/dominated omissions; inverse-image optimal controls when onto.\\
Missing middle & Four compatible optimal composite pairs, two with a
middle policy.\\
Concurrency and horizon & The capacity choice reverses with the coupling;
two-step preparation gives \(2\), rolling one-step evaluation gives \(0\).\\
\bottomrule
\end{tabularx}
\end{center}

The main arena emits a fair bit under root control A; control B retains
both bit witnesses with law \((1,0)\) and pays \(1/2\). At each ensuing
history, preparation costs \(1/2\) and increases the final work reward from
\(1\) to \(2\). Raw bit \(1\) adds \(1/4\) to the work reward. At the last
stage, work earns that reward and rest earns zero. Terminal values are
zero. These data give root action values \(13/8\) and \(2\).
Direct path evaluation does not use the backward values. It separately
computes occupancy and checks \eqref{eq:loss}.

The finite enumerations do not prove the real-score theorem, arbitrary
map naturality, or the general category laws; their written proofs appear
above. The associativity diagnostic checks \(25\) rational terminal-score
pairs; Proposition~\ref{prop:naive} supplies the all-real-scores argument.
No empirical performance or proof-assistant certificate is asserted.

\section{Source boundary}
Repository references below are pinned to
\texttt{777a28ab\ldots}, whose full identifier is recorded in the
companion source manifest. The contextual-composition files used here
were first read at \texttt{f325c370\ldots}; comparison with the pinned
revision found no changes in that project. Source results are credited
at their stated hypotheses. The pilot does not rerun historical repository
proof-assistant checks or promote a source's conditional native theorem
to an unconditional realization theorem.

\begin{thebibliography}{9}
\bibitem{contextads}
M.~Capucci and D.~J.~Myers.
\emph{Contextful Arrows via Wreath Products}.
Version 2, 2026; first version, 2024.
\href{https://arxiv.org/abs/2410.21889v2}{arXiv:2410.21889v2}.

\bibitem{terminology}
Contextual Composition Programme.
\emph{Interactads: a provisional terminology note}.
\reposrc{foundations/terminology/Interactad_Terminology_Note.tex}{Pinned TeX source}.
The package to be named; Two qualified usages.

\bibitem{relative}
Contextual Composition Programme.
\emph{One relative construction for context content and effects}.
\reposrc{foundations/relative-construction/delivery/One_relative_construction_for_context_content_and_effects.tex}{Pinned TeX source}.
In particular, \emph{Comparison maps and controlled forgetting}.

\bibitem{conditional}
Contextual Composition Programme.
\emph{Full-control-history conditional families of open timeline amalgamations}.
\reposrc{foundations/native-whole-scope-feedback/Conditional_Open_Timeline_Families.tex}{Pinned TeX source}.
\emph{Unselected law-labelled control histories} and
\emph{Exact relative extension and descent theorem}; theorems
\emph{Unselected conditional-family normal form} and
\emph{Relative lift fibre, including uncovered raw histories}.

\bibitem{transport}
Contextual Composition Programme.
\emph{Policy relations compose only laxly: the missing-middle obstruction}.
\reposrc{foundations/native-whole-scope-feedback/Policy_Relational_Transport.tex}{Pinned TeX source}.
Theorems \emph{Normal lax relational transport} and
\emph{All-policy lift of presented open assembly}; the separate strict witness.

\bibitem{nonblocking}
Contextual Composition Programme.
\emph{Nonblocking recurrent interaction: history-law implementations and certified opaque exports}.
\reposrc{foundations/native-whole-scope-feedback/Nonblocking_History_Controllers_and_Certified_Exports.tex}{Pinned TeX source}.
\emph{Global admissibility and actual implementations};
\emph{Calibrated complete-plan gluing} (C1--C2);
\emph{Relative native finality with implementations}.

\bibitem{cybernetics}
M.~Capucci, B.~Gavranovi\'c, J.~Hedges, and E.~F.~Rischel.
\emph{Towards Foundations of Categorical Cybernetics}.
Version 2, 2022.
\href{https://arxiv.org/abs/2105.06332v2}{arXiv:2105.06332v2}.

\bibitem{rl}
J.~Hedges and R.~Rodr\'iguez Sakamoto.
\emph{Reinforcement Learning in Categorical Cybernetics}.
Version 2, 2025.
\href{https://arxiv.org/abs/2404.02688v2}{arXiv:2404.02688v2}.
\end{thebibliography}
\end{document}
